%!TEX root = ../thesis_main.tex
%!TEX encoding = UTF-8 Unicode

\section{Experimental set-up}

The experimental gasifier is a pilot two-stage downdraft gasifier based on the NOTAR{\copyright} technology, extensively described by \citet{Berger2018}. This gasifier design is structured in three main sections (Figure \ref{fig:notar}):
\begin{enumerate}
\item the pyrolysis zone, where biomass is dried and converted into char and volatile compounds by an oxidative pyrolysis;
\item the oxidation zone, where volatiles are oxidized and cracked in homogeneous reactions at high temperature;
\item the reduction zone, where heterogeneous char gasification reactions convert the gaseous products into useful \ce{CO} and \ce{H2}.
\end{enumerate}

Air injection is split into two stages: primary air at the top of the pyrolysis zone and secondary air in the oxidation zone. At both stages, air is injected through four nozzles evenly distributed around the gasifier. The independent control of the pyrolysis, and the high-temperature homogeneous oxidation zone where tars are cracked and oxidized by the secondary air injection, are the key features enabling a very low tar content below 100\,mg/Nm$^3$.

The nominal biomass consumption is 50\,kg/h, corresponding to 250\,kW of biomass LHV and up to 175\,kW of syngas LHV. Wood chips are loaded in a hopper from which they are fed to the gasifier using a worm screw, via an airlock to prevent any escape of syngas. Temperature is measured at 20 points located in spiral along the gasifier height: 11 in the pyrolysis zone, 3 in the oxidation zone and 8 in the reduction zone. At the bottom of the gasifier, the reduction bed lies on an ash rack which can be actuated to help dislocate the bed and extract the ashes when pressures losses build up.

\begin{figure}[!ht]
	\centering
	\includegraphics[width=\textwidth]{notar_design.pdf}
	\caption{The two-stage gasifier design.\newline Pyrolysis volatiles are oxidized and cracked in the hot gas-phase zone}
	\label{fig:notar_design}
\end{figure}


\begin{figure}[!ht]
	\centering
	\includegraphics[width=0.8\textwidth]{notar.pdf}
	\caption{The two-stage gasifier design.\newline Pyrolysis volatiles are oxidized and cracked in the hot gas-phase zone \cite{Berger2018}}
	\label{fig:notar}
\end{figure}

Syngas cleaning is required to protect any downstream piping or equipment from clogging due to the condensation of remaining tar compounds. Cooling and cleaning of raw syngas is done in four main steps at the gasifier outlet (Figure \nolinebreak\ref{fig:cleaning}). Fine ash in suspension are first removed by a dust cyclone (4). Syngas is then mixed with water in a Venturi scrubber right upon entering a counter-current heat exchanger (5) where water and tar condense out of the cooled syngas stream. The following wet cyclone (6) ensures that fine droplets of water entrained by the flow are removed, and the granular filter (7) acts as a final trap for remaining solid particles or droplets. The clean syngas flow is then measured with an orifice flow meter, before flaring or use.

\begin{figure}[!ht]
	\centering
	\includegraphics[width=\textwidth]{cleaning.pdf}
	\caption{Gasifier and syngas cleaning: (1-2) ash tank and extraction, (3) gasifier, (4) cyclone and fine ash collector, (5) condenser, (6) wet cyclone, (7) granular filter, (8) flare and (9) water decantation tank \cite{Berger2018}.}
	\label{fig:cleaning}
\end{figure}

The steam injection system consists of an electric steam generator of 22 kW producing up to 25\,kg/h of saturated steam, followed by a steam superheater. Steam flow regulation is done with an orifice flow meter and a control valve. Insulation and trace heating are used on the steam lines to maintain the steam temperature until the injection nozzles to the oxidation zone.

Oxygen is produced with a Pressure-Swing Adsorption (PSA) oxygen generator, with an oxygen purity of $x_{\ce{O2},\text{oxygen}}=95\%$. Steam, air and oxygen are injected at the secondary stage through 4 injection nozzles spread around its circumference. Air and steam are mixed inside the nozzles, ensuring their homogeneous mixing before injection, while oxygen is injected through concentric nozzles for safety reasons.

Syngas composition after cleaning is determined using an on-line gas analyser. $\ce{CO}$, $\ce{H2}$ and $\ce{CH4}$ are quantified with a Siemens ULTRAMAT 23 non-dispersive infrared sensor (NDIR), which also contains a paramagnetic sensor used to detect the presence of $\ce{O2}$. A thermal conductivity detector (TCD) Siemens CALOMAT 62 is coupled to measure the $\ce{H2}$ concentration. $\ce{N2}$ is calculated by difference as $100\%-[\ce{CO}]-[\ce{CO2}]-[\ce{CH4}]-[\ce{H2}]$, and actually includes all other inert species like \ce{Ar}.

Trace compounds are also detected and some of them quantified by gas chromatography using an on-line GC-MS-FID, where the mass spectrometer (MS) and flame ionization detector (FID) are installed in parallel. The six dominant species identified and quantified are benzene, ethylene, ethane, methanol, propene and acetone. The list of all light hydrocarbons found as traces in the syngas is given in Table \ref{tab:GCMS_lightHCs} in the Appendix. 


\section{Experimental conditions}

\begin{figure}[!ht]
    \centering
    \includegraphics[width=0.5\textwidth]{woodchips.pdf}
	\caption{Oven-dried wood chips of 10-40\,mm are used.}
	\label{fig:woodchips}
\end{figure}

Wood chips of 10-40\,mm with less than 10\% moisture are used as feedstock for all experiments. Two different batches were used over the experimental campaigns, with very similar properties. Their detailed composition is given in Table \ref{tab:woodchips} in the Appendix.

Experimental conditions are detailed in Table \ref{tab:xpplan}. The two primary factors of interest are the Steam-to-Oxygen molar ratio $\mathrm{SOR}$ (Eq. (\ref{eq:SOR})) and the equivalent air enrichment $x_{\ce{O2}}$ [\%] (Eq. (\ref{eq:xO2})). The $\mathrm{SOR}$ is varied from 0 to 2 and represents the steam flow normalized by the total input of \ce{O2} from air and oxygen flows. Oxygen injection at the secondary stage results in an increase of $x_{\ce{O2}}$, ranging from 21\% with air to 47\% when secondary air is replaced by oxygen. Intermediate values correspond to a mixture of secondary air and oxygen.

Secondary control factors are the total equivalent air flow $Q_\text{air,eq}$ [Nm$^3$/h] (Eq. (\ref{eq:Qair})), the secondary-to-primary \ce{O2} ratio $\beta$ (Eq. (\ref{eq:AR})) and the injection temperature of steam $T_{st}$. $Q_\text{air,eq}$ corresponds to the air flow containing an equivalent amount of \ce{O2} as the combined air and oxygen flows, and should primarily control the gasifier power. It was varied between 60 and 70\,Nm$^3$/h. $\beta$ represents the ratio of secondary and primary \ce{O2} flow from the combined air and oxygen flows. It is a control variable that must be adjusted during operation to maintain the pyrolysis front stability, and is usually between 2.0 and 2.5. Finally, steam was either saturated (LT, 100-150°C) or superheated (HT, 250-500°C).

\begin{eqnarray}
\label{eq:SOR} \mathrm{SOR} &=& \frac{\dot n_\mathrm{H2O_g}   }{ \dot n_{\ce{O2}  }} \;\; = \;\;  \frac{  \dot m_\text{steam}/M_\text{steam}   }{  x_{\ce{O2},\text{air}} Q_\text{air,eq} /V_0}\\
\label{eq:xO2} x_{\ce{O2}} &=& \frac{  x_{\ce{O2},\text{air}}Q_\text{air,primary} + x_{\ce{O2},\text{air}}Q_\text{air,secondary} + x_{\ce{O2},\text{oxygen}}Q_\text{oxygen,secondary}                }{       Q_\text{air,primary} + Q_\text{air,secondary} + Q_\text{oxygen,secondary}           }\\
\label{eq:Qair} Q_\text{air,eq} &=& Q_\text{air,primary} + Q_\text{air,secondary} + Q_\text{oxygen,secondary}\frac{x_{\ce{O2},\text{oxygen}}}{x_{\ce{O2},\text{air}}}\\
\label{eq:AR} \beta &=& \frac{   x_{\ce{O2},\text{air}}Q_\text{air,secondary} + x_{\ce{O2},\text{oxygen}}Q_\text{oxygen,secondary}   }{      x_{\ce{O2},\text{air}}Q_\text{air,primary}        } 
\end{eqnarray}

\renewcommand\theadfont{}

\begin{table}[!ht]
\renewcommand\theadalign{c}
\centering
\begin{tabular}{llccccccS[table-number-alignment=right, table-figures-integer = 3]}
\thead[l]{\nth{2} stage} & \thead[l]{Run ID} &  \thead{$\mathrm{SOR}$\\\;}  & \thead{$x_{\ce{O2}}$\\\%} &  \thead{$Q_\text{air,eq}$\\Nm$^3$/h}  &  \thead{$\beta$\\\;} &     \thead{$T_\mathit{st}$\\°C}    & \thead{Wood\\batch} & \multicolumn{1}{c}{\thead{Duration\\min}}\\
\specialrule{0pt}{2pt}{2pt}
\multirow{4}{*}{Air} & A1 &  0.00 &  21.0 &     59.4 &  2.26 &        &      1 &        79 \\
                               & A2 &  0.00 &  21.0 &     58.6 &  1.97 &        &      2 &        74 \\
                               & A3 &  0.00 &  21.0 &     68.8 &  2.49 &        &      2 &       108 \\
                               & A4  &  0.00 &  21.0 &     69.3 &  2.99 &        &      2 &        90 \\
\specialrule{0pt}{2pt}{2pt}
\multirow{3}{*}{Air-steam} & B1 &  0.86 &  21.0 &     69.2 &  1.94 &    104 &      1 &        58 \\
                               & B2 &  2.00 &  21.0 &     59.5 &  2.00 &    103 &      1 &        99 \\
                               & B3 &  2.03 &  21.0 &     58.8 &  2.00 &    271 &      1 &        67 \\
\specialrule{0pt}{2pt}{2pt}
\multirow{2}{*}{Air-\ce{O2}} & C1 &  0.00 &  28.2 &     64.8 &  1.63 &        &      2 &       108 \\
                               & C2 &  0.00 &  32.3 &     70.0 &  2.55 &        &      2 &        70 \\
\specialrule{0pt}{2pt}{2pt}
\multirow{2}{*}{\ce{O2}} & D1 &  0.00 &  47.3 &     59.8 &  2.53 &        &      2 &       100 \\
                               & D2 &  0.00 &  47.5 &     70.0 &  2.55 &        &      2 &        78 \\
\specialrule{0pt}{2pt}{2pt}
\multirow{6}{*}{\ce{O2}-steam} & E1 &  1.09 &  44.0 &     60.6 &  2.07 &    117 &      2 &        72 \\
                               & E2 &  0.99 &  47.2 &     59.8 &  2.52 &    146 &      2 &        68 \\
                               & E3 &  0.99 &  43.8 &     60.2 &  2.04 &    399 &      2 &        82 \\
                               & E4 &  1.98 &  47.3 &     60.1 &  2.53 &    127 &      2 &        97 \\
                               & E5 &  1.97 &  47.5 &     60.3 &  2.55 &    460 &      2 &        67 \\
                               & E6 &  1.97 &  47.4 &     70.3 &  2.53 &    134 &      2 &        91 \\
\end{tabular}
\caption{Experimental conditions of the 17 experimental points.\label{tab:xpplan}}
\end{table}


%\begin{table}[!ht]
%\renewcommand\theadalign{c}
%\centering
%\begin{tabular}{llccccccS[table-number-alignment=right, table-figures-integer = 3]}
%\thead[l]{\nth{2} stage} & \thead[l]{Run ID} &  \thead{$\mathrm{SOR}$\\\;}  & \thead{$x_{\ce{O2}}$\\\%} &  \thead{$Q_\text{air,eq}$\\Nm$^3$/h}  &  \thead{$\beta$\\\;} &     \thead{$T_\text{steam}$\\°C}    & \thead{Wood\\batch} & \multicolumn{1}{c}{\thead{Duration\\min}}\\
%\specialrule{0pt}{2pt}{2pt}
%\multirow{4}{*}{Air} & air60\_A &  0.00 &  21.0 &     59.4 &  2.26 &        &      1 &        79 \\
%                               & air60\_B &  0.00 &  21.0 &     58.6 &  1.97 &        &      2 &        74 \\
%                               & air70\_A &  0.00 &  21.0 &     68.8 &  2.49 &        &      2 &       108 \\
%                               & air70\_B &  0.00 &  21.0 &     69.3 &  2.99 &        &      2 &        90 \\
%\specialrule{0pt}{2pt}{2pt}
%\multirow{3}{*}{Air-steam} & air70\_st10LT &  0.86 &  21.0 &     69.2 &  1.94 &    104 &      1 &        58 \\
%                               & air60\_st20LT &  2.00 &  21.0 &     59.5 &  2.00 &    103 &      1 &        99 \\
%                               & air60\_st20HT &  2.03 &  21.0 &     58.8 &  2.00 &    271 &      1 &        67 \\
%\specialrule{0pt}{2pt}{2pt}
%\multirow{2}{*}{Air-\ce{O2}} & air\ce{O2}65 &  0.00 &  28.2 &     64.8 &  1.63 &        &      2 &       108 \\
%                               & air\ce{O2}70 &  0.00 &  32.3 &     70.0 &  2.55 &        &      2 &        70 \\
%\specialrule{0pt}{2pt}{2pt}
%\multirow{2}{*}{\ce{O2}} & \ce{O2}60 &  0.00 &  47.3 &     59.8 &  2.53 &        &      2 &       100 \\
%                               & \ce{O2}70 &  0.00 &  47.5 &     70.0 &  2.55 &        &      2 &        78 \\
%\specialrule{0pt}{2pt}{2pt}
%\multirow{6}{*}{\ce{O2}-steam} & \ce{O2}60\_st10LT\_A &  1.09 &  44.0 &     60.6 &  2.07 &    117 &      2 &        72 \\
%                               & \ce{O2}60\_st10LT\_B &  0.99 &  47.2 &     59.8 &  2.52 &    146 &      2 &        68 \\
%                               & \ce{O2}60\_st10HT &  0.99 &  43.8 &     60.2 &  2.04 &    399 &      2 &        82 \\
%                               & \ce{O2}60\_st20LT &  1.98 &  47.3 &     60.1 &  2.53 &    127 &      2 &        97 \\
%                               & \ce{O2}60\_st20HT &  1.97 &  47.5 &     60.3 &  2.55 &    460 &      2 &        67 \\
%                               & \ce{O2}70\_st20LT &  1.97 &  47.4 &     70.3 &  2.53 &    134 &      2 &        91 \\
%\end{tabular}
%\caption{Experimental conditions of the 17 experimental points.\label{tab:xpplan}}
%\end{table}

\section{Response variables and steady-state results selection}

The primary response variables are the syngas composition (\ce{CO}, \ce{H2}, \ce{CO2}, \ce{CH4}) and $\mathrm{LHV}$ on dry basis, after cooling (Eq. (\ref{eq:LHV})). The syngas \ce{H2}/\ce{CO} ratio is observed, as it is strongly influenced by the injection of steam. The Energy Density $\mathrm{ED}$ (Eq. (\ref{eq:ED})) and Adiabatic Flame Temperature $\mathrm{AFT}$ (Eq. (\ref{eq:AFT})) of the stoichiometric air-syngas mixture are computed as secondary responses. The $\mathrm{ED}$ controls the volumetric power of an Internal Combustion Engine (ICE), relevant for CHP applications. The $\mathrm{AFT}$ gives an idea of the flame temperature that can be reached, useful for high-temperature burners for industrial applications. In Equation (\ref{eq:AFT}), $n_i$ and $h_{s,i}$ are the amount [kmol] and sensible enthalpy [kJ/kmol] of each product of the complete stoichiometric combustion of one kmol of syngas, and $V_0=\frac{RT_0}{p_0}=22.414$\,Nm$^3$/kmol is the normal volume of 1\,kmol under NTP conditions of 0°C and 1\,bar.

The gasifier power is observed, both in terms of biomass consumption $P_\text{wood}$ and syngas production $P_\text{syngas}$. The air-fuel Equivalence Ratio (ER) is computed (Eq. (\ref{eq:ER_XP})), as well as the Cold Gas Efficiency $\mathrm{CGE}$ (Eq. (\ref{eq:CGE_XP})) and the Carbon Conversion Efficiency $\mathrm{CCE}$ (Eq. (\ref{eq:CCE})). As the injection of oxygen can result in a problematic temperature increase, the average temperature in the oxidation zone $T_\text{oxi,avg}$ is also monitored. Finally, the amount and composition of tars and lighter hydrocarbons are of primary interest in this study.

\begin{eqnarray}
\label{eq:LHV} \mathrm{LHV} &=& x_{\ce{CO}}\mathrm{LHV}_{\ce{CO}} + x_{\ce{H2}}\mathrm{LHV}_{\ce{H2}} + x_{\ce{CH4}}\mathrm{LHV}_{\ce{CH4}}\\
\label{eq:ED} \mathrm{ED} &=& \left.\mathrm{LHV}\right/ \left(1 + \frac{0.5x_{\ce{CO}} + 0.5x_{\ce{H2}} + 2x_{\ce{CH4}}}{x_{\ce{O2},\text{air}}}\right)\\
\label{eq:AFT} \mathrm{AFT} & \equiv &\mathrm{LHV} \;\;= \;\; \sum_{i\in P} n_i h_{s,i}(\mathrm{AFT})/V_0\\
\label{eq:ER_XP} \mathrm{ER} &=& \frac{x_{\ce{O2}}Q_\text{air,eq}/V_0}{\left(1+\frac y4 - \frac x2\right)\dot m_{\text{biomass},\mathit{daf}}/M_{\ce{CH_yO_xN_z}}}\\
\label{eq:CGE_XP} \mathrm{CGE}&=&\frac{P_\text{syngas}}{P_\text{biomass,\textit{wb}}}\;\; = \;\;\frac{Q_\text{syngas}\mathrm{LHV}}{\dot m_{\text{biomass},wb}\mathrm{LHV}_{\text{biomass},wb}}\\
\label{eq:CCE} \mathrm{CCE} &=& \frac{\dot n_{\ce{C},\text{syngas}}}{\dot n_{\ce{C},\text{biomass}}} =  \frac{\left(x_{\ce{CO}}+x_{\ce{CO2}}+x_{\ce{CH4}}\right) Q_\text{syngas}/V_0}{\dot m_{\text{biomass},\mathit{daf}}/M_{\ce{CHyO_xN_z}}}
\end{eqnarray}

Due to its scale and substantial mass of concrete, the gasifier has a big inertia and several hours of operation are required to reach stable steady-state operating points. A typical evolution of the composition over a day is shown in Figure \ref{fig:heating}. Experimental results presented in this paper are based on stable operation periods of a duration varying between 58 and 108 minutes. These periods were selected based on the stability of the syngas composition and $\mathrm{LHV}$.

Syngas composition fluctuates with a regular pattern, illustrated in Figure \ref{fig:fluctuations}, due to the non-continuous feeding of charcoal from the pyrolysis zone to the reduction zone. The exact stability periods were adapted so that their duration matches an integer number of fluctuations, to avoid biasing the average value. Based on the selected periods, the stability of the syngas $\mathrm{LHV}$ was assessed more precisely by computing its relative average slope $\frac{1}{\mathrm{LHV}}\frac{d\mathrm{LHV}}{dt}$ using a linear regression. This value was used to detect and discard unsteady experimental points, defined by a $\mathrm{LHV}$ increasing at a rate of more than 5\% per hour.

Experimental results are averaged over each considered stable period. These values are then used to fit regression models including second-order and interaction terms on $x_{\ce{O2}}$ and $\mathrm{SOR}$, and first-order terms on $Q_{\text{air,eq}}$ and $\beta$. The steam temperature $T_{st}$ was not found to have any statistically relevant impact and was discarded from the models. The obtained regression models are used to synthesize experimental results and observe trends.

Wood chips consumption $\dot m_{\text{wood},\mathit{wb}}$ is computed based on the time between successive airlock loadings. The mass of wood chips in the airlock is determined by dividing the total mass loaded to the hopper by the corresponding number of airlock loadings. For each of the two batches of wood chips, the resulting mass of chips per airlock is averaged over all the associated hopper loadings, giving an individual value for each biomass batch. The "instantaneous" biomass consumption is then computed by dividing the mass of chips per airlock by the duration of each airlock loading. This value is smoothed out by a "rolling average" with a triangular window of 60 minutes, which is then averaged over the stable operation period to obtain the "local average", used as reference for the given run (Figure \ref{fig:mbois}).

%The rate of biomass consumption is estimated based on the mass of wood chips per hopper loading, and the time difference between successive loadings. The mass of wood chips per load is determined using a 1:1 model of the hopper on a scale, filled manually.  Based on the biomass, air and steam flow, the wood input power, air-to-fuel Equivalence Ratio (ER) and the Steam-to-Biomass mass Ratio (SBR) corresponding to each run are given in Table \ref{tab:SBR_ER_Pbois}. While the wood consumption is proportional to the total air flow, it is not clearly affected by steam injection.
%As the duration of a single load is around 20 minutes, the uncertainty on the actual instantaneous biomass consumption is high and difficult to assess. 

%Based on equilibrium model simulations, the optimal Equivalence Ratio is expected to decrease as steam gasification partly replaces air gasification. However, the ER is not controlled but results from the rate of biomass consumption, depending on the rate of consumption of the char bed. Practically, steam injection can be expected to either accelerate the char reduction by reacting with it, or slow it down by inducing a local temperature drop. These two effects experimentally appear to partly offset each other, yielding a very small increase of the ER, for an approximately constant value around 0.37. This average value is used in equilibrium computations below to compare experimental results to simulations.

\begin{figure}[!ht]
	\centering
	\includegraphics[width=0.7\textwidth]{TGPheating.pdf}
	\caption{A long heating time is required to reach stable operation.\newline\textit{Run IDs:} air60{\_}B \textit{and} air\ce{O2}65.}
	\label{fig:heating}
\end{figure}

\begin{figure}[!ht]
	\centering
	\includegraphics[width=0.7\textwidth]{fluctuations.pdf}
	\caption{Syngas composition and LHV fluctuate with a period of 15-25\,minutes.\newline\textit{Run ID:} \ce{O2}60{\_}st20LT}
	\label{fig:fluctuations}
\end{figure}

\begin{figure}[!ht]
	\centering
	\includegraphics[width=0.7\textwidth]{mbois.pdf}
	\caption{Wood chips consumption is computed based on the duration of individual airlocks, averaged out first with a rolling average then a local average to consider border effects.\newline\textit{Run ID:} \ce{O2}60{\_}st20LT\label{fig:mbois}}
\end{figure}

%\begin{figure}[!ht]
%	\centering
%	\includegraphics[width=0.6\textwidth]{xpplan.png}
%	\caption{Experimental conditions. LT and HT refer to Low and High Temperature steam. LT steam is injected at saturation temperature (100°C) while HT steam is superheated (250°C).}
%	\label{fig:xpplan}
%\end{figure}


\section{Experimental results}

Experimental results and the associated regression models are presented and discussed in this section, and some of them compared with the equilibrium model predictions. First, the syngas composition, \ce{H2}/\ce{CO} ratio, LHV, ED and AFT are observed. The gasifier power, CGE and CCE are analyzed, followed by the impact of steam and oxygen on the temperature in the oxidation zone. Tars and light hydrocarbons are finally discussed. The complete results of syngas composition, LHV, biomass and syngas flow and power, and efficiency, are provided in Tables \ref{tab:xpres} and \ref{tab:xpres_eff} in the Appendix.

In all figures and tables presented in this section, regression-based results are computed using $Q_\text{air,eq}=65$ and $\beta=2.25$ if not otherwise specified. In the graphs, the following code applies to the markers used to represent experimental results: $Q_\text{air,eq}$=60 (circle), 65 (diamond) or 70\,Nm$^3$/h (square); $x_{\ce{O2}}$=21\% (empty), 32\% (empty with a dot) or 47\% (full); $\mathrm{SOR}$=0 (dark red), 1 (medium red) or 2 (light red).


\subsection{Syngas composition and \ce{H2}/\ce{CO} ratio}
\label{sec:comp}

The volume composition of dry syngas is given in Table \ref{tab:comp}. The concentrations of \ce{H2}, \ce{CO}, \ce{CO2} and \ce{N2} are shown in Figures \ref{fig:comp_air} and \ref{fig:comp_O2} for different values of the SOR, respectively with air and oxygen at the secondary stage. These values are computed using the regression model fitted on the experimental data, including $\nth{2}$ order and interaction terms for $\mathrm{SOR}$ and $x_\mathrm{\ce{O2}}$, and $\nth{1}$ order terms for $Q_\text{air,eq}$ and $\beta$.

Steam injection boosts the production of \ce{H2} and \ce{CO2} at the expense of \ce{CO}. This effect is due to the Water-Gas Shift reaction (Eq. (\ref{eq:WGS})) being shifted towards the products, according to Le Chatelier's principle. The use of oxygen, instead of air, reduces the concentration of \ce{N2} in the syngas, scaling up \ce{H2}, \ce{CO}, \ce{CO2} and \ce{CH4} accordingly.

\renewcommand\theadfont{}
\renewcommand\theadalign{c}
\begin{table}[!ht]
\centering
\begin{tabular}{lcc@{\hspace{0.5cm}}ccccc}
 \thead[l]{\;\\\nth{2} stage} & \thead{$x_{\ce{O2}}$\\\%} & \thead{SOR\\\;} & \thead{\ce{CO}\\\%} & \thead{\ce{H2}\\\%} & \thead{\ce{CH4}\\\%} & \thead{\ce{CO2}\\\%} & \thead{\ce{N2}\\\%}\\
\multirow{3}{*}{Air} & \multirow{3}{*}{21} & 0 &            18.4 &            16.8 &             0.73 &             12.2 &            51.9 \\
                         &                     & 1 &            15.2 &            21.1 &             0.90 &             14.6 &            48.1 \\
                         &                     & 2 &            12.4 &            23.6 &             1.00 &             17.2 &            45.7 \\
\specialrule{0pt}{2pt}{2pt}
Air-\ce{O2} & 32 & 0 &            24.2 &            24.5 &             1.21 &             14.9 &            35.1 \\
\specialrule{0pt}{2pt}{2pt}
\multirow{3}{*}{\ce{O2}} & \multirow{3}{*}{47} & 0 &            29.1 &            29.2 &             1.21 &             18.6 &            21.9 \\
                         &                     & 1 &            23.5 &            33.9 &             1.42 &             21.7 &            19.4 \\
                         &                     & 2 &            18.4 &            36.8 &             1.54 &             24.9 &            18.3 \\        
		\end{tabular}
\caption{Syngas composition, predicted by the regression model built on experimental results. $Q_\text{air,eq}=65$, $\beta=2.25$\label{tab:comp}}
\end{table}


\begin{figure}[!ht]
%	\begin{minipage}{0.49\textwidth}
		\centering
		\includegraphics[width=0.6\textwidth]{comp_21.pdf}
		\caption{Dry syngas composition with air \textit{(regression model with $Q_\text{air,eq}=65$ and $\beta=2.25$)}.}
		\label{fig:comp_air}
%	\end{minipage}
\end{figure}
%	\hfill
\begin{figure}[!ht]
%	\begin{minipage}{0.49\textwidth}
		\centering
		\includegraphics[width=0.6\textwidth]{comp_47.pdf}
		\caption{Dry syngas composition with oxygen \textit{(regression model with $Q_\text{air,eq}=65$ and $\beta=2.25$)}.}
		\label{fig:comp_O2}
%	\end{minipage}
\end{figure}


The shift of syngas from \ce{CO} towards \ce{H2} with steam injection translates into a linear increase of the \ce{H2}/\ce{CO} ratio: from almost 1.0 without steam up to 2.0 at $\mathrm{SOR}=2$ (Figure \ref{fig:H2_CO}). Experimental results are fitted by a regression model based only on $\mathrm{SOR}$ and $x_{\ce{O2}}$, as there is no statistically valid impact of the other factors. The \ce{H2}/\ce{CO} ratio is a relevant metric for some applications: a high value is desirable for the synthesis of advanced biofuels such as methane or methanol, or for \ce{H2} production. Moreover, the combustion behaviors of \ce{H2} and \ce{CO} are different, and their relative share can affect the flame properties such as the laminar flame speed and ignition temperature \cite{Sung2008}.

\begin{figure}[!ht]
	\centering
	\includegraphics[width=0.6\textwidth]{H2_CO.pdf}
	\caption{The \ce{H2}/\ce{CO} ratio increases linearly with the $\mathrm{SOR}$, doubling from 1.0 to 2.0. \textit{Regression (lines) and experimental points (markers).}}
	\label{fig:H2_CO}
\end{figure}

%\begin{figure}[!ht]
%	\centering
%	\includegraphics[width=0.7\textwidth]{H2_CO_EM.pdf}
%	\caption{\ce{H2}/\ce{CO} ratio\\ predicted by the equilibrium model\textbf{caption to complete}\label{fig:H2_CO_EM}}
%\end{figure}



Beyond these primary observations of the dry syngas composition, the impacts of oxygen and steam can be better observed, understood and compared with the equilibrium model by avoiding the effect of the variable syngas dilution. The normalized composition $x_{i,\mathrm{norm}}$ [\%] is calculated by dividing the volume concentration by the total amount of carbon in the syngas (Eq. (\ref{eq:x_norm})). The amount of carbon is representative of the amount of biomass converted into syngas, disregarding the part of carbon lost in the ashes.

\begin{equation}\label{eq:x_norm}
x_{i,\mathrm{norm}} = \frac{x_i}{x_{\ce{CO}}+x_{\ce{CO2}}+x_{\ce{CH4}}}
\end{equation}

With oxygen, more \ce{H2} and \ce{CO} and less \ce{CO2} are produced than with air, in normalized composition (Figure \ref{fig:comp_norm}). The lower dilution reduces the amount of energy lost as syngas sensible enthalpy, and less complete combustion to \ce{CO2} and \ce{H2O} is therefore required to sustain the process energy needs. This results in a lower Equivalence Ratio, and a higher production of \ce{CO} and/or \ce{H2}. 

%When replacing air by oxygen, the equilibrium model mainly predicts an increase in \ce{CO} production, not affecting \ce{H2} (Figure \ref{fig:comp_norm_EM}). The experimental results show the opposite trend: most gain is on \ce{H2} rather than \ce{CO} production. Based on the model, the equilibrium temperature is expected to be higher with oxygen, pushing the water-gas shift reaction (Eq. (\ref{eq:WGS})) towards more \ce{CO} and less \ce{H2}. Practically, no substantial difference on the syngas outlet temperature is observed between air and oxygen results, which can explain why this theoretical effect is not matched experimentally. The difference can also be partly attributed to variations in the biomass composition: a slightly more humid and hydrogenated biomass was used during the experiments with \ce{O2}, with a global H/C ratio of 1.58 instead of 1.45 (Table \ref{tab:woodchips}), thus favoring \ce{H2} production. The impact of steam on the normalized composition is well predicted by the equilibrium model, except for the slight quadratic trend that is not observed experimentally.


\begin{figure}[!ht]
%	\begin{minipage}{0.49\textwidth}
		\centering
		\includegraphics[width=0.6\textwidth]{comp_norm.pdf}
		\caption{Syngas composition\newline normalized by the amount of carbon:\newline experimental results and regressions.}
		\label{fig:comp_norm}
%	\end{minipage}
\end{figure}
%	\hfill
%\begin{figure}[!ht]
%	\begin{minipage}{0.49\textwidth}
%		\centering
%		\includegraphics[width=0.6\textwidth]{comp_norm_EM.pdf}
%		\caption{Syngas composition\newline normalized by the amount of carbon:\newline equilibrium model predictions.}
%		\label{fig:comp_norm_EM}
%	\end{minipage}
%\end{figure}


\subsection{LHV, Energy Density and Adiabatic Flame Temperature}

The volume-based LHV of dry syngas is the best metric to characterize its intrinsic energy content. It is preferable to HHV as in most applications of syngas, the latent enthalpy of vaporisation of water is not valorized. The use of HHV would artificially favor a \ce{H2}-rich syngas, as \ce{H2} has a higher HHV but lower LHV than \ce{CO}. The use of a volume basis (MJ/Nm$^3$) is more relevant than a mass basis (MJ/kg): as syngas is a gaseous fuel, its storage, transport and control are mainly constrained by its volume rather than weight. The use of a mass basis would also artificially favor a \ce{H2}-rich syngas, due to its low density, while mass is not a real issue when dealing with gaseous fuels.

Thanks to the reduced dilution and higher production of \ce{H2} and \ce{CO}, syngas LHV is increased up to $+65\%$ when replacing air by oxygen at the secondary stage: from 4.40 to 7.26\,MJ/Nm$^3$ (Figure \ref{fig:LHV} and Table \ref{tab:LHVEDAFT_species}). The effect of air enrichment is quadratic: LHV is already 6.14\,MJ/Nm$^3$ at $x_{\ce{O2}}=32$\% ($+40\%$), which corresponds to a partial secondary air enrichment of around 40\% \ce{O2}. The equilibrium model prediction (Figure \ref{fig:LHV_EM}) is in good agreement with experimental results, although the quadratic effect is milder: from 4.50 with air to 7.33\,MJ/Nm$^3$ with \ce{O2}, and 6.02 at $x_{\ce{O2}}=32$\% ($+34\%$ compared to air).  Due to the limited number of experimental results in the Air-\ce{O2} conditions, the trend predicted by the equilibrium model is probably more representative than the regression model; the p-Value of $\mathrm{SOR}^2$ term in the regression model is 0.45, hence not statistically valid.

Steam seems to have a slightly positive impact on the LHV with air (up to around 4.50\,MJ/Nm$^3$), but this very small effect is not statistically representative. With oxygen, steam reduces the syngas LHV, especially at SOR=2 where it drops to 6.85\,MJ/Nm$^3$ ($-5\%$). Compared to the reference air conditions, the LHV improvement with oxygen and SOR=2 is +56\%. This effect of steam is explained by the higher amount \ce{CO2} produced by the water-gas-shift reaction (Eq. (\ref{eq:WGS})), which dilute the syngas.


Based on the equilibrium model, the negative impact of steam is expected to be much stronger, lowering the LHV to 6.32\,MJ/Nm$^3$ with \ce{O2} and to 4.03\,MJ/Nm$^3$ with air. As explained in Section \ref{sec:eqmodel}, the Equivalence Ratio (ER) is assumed not to vary with steam injection, to avoid unrealistic results. However, experimental results show that the ER is actually lower when injecting steam, especially in the runs with air: steam partly replaces \ce{O2} as gasification agent, although to a much lower extent than predicted by the equilibrium model with CGE optimization.


\begin{figure}[!ht]
%	\begin{minipage}{0.49\textwidth}
		\centering
		\includegraphics[width=0.6\textwidth]{LHV.pdf}
		\caption{Syngas LHV is increased to 7.3\,MJ/Nm$^3$ ($+65\%$) when replacing secondary air by oxygen,\newline lowered to 6.9\,MJ/Nm$^3$ ($+56\%$) with $\mathrm{SOR}=2$. \textit{Lines: regression model ($Q_\text{air,eq}=65$, $\beta=2.25$). Markers: experimental points.}}
		\label{fig:LHV}
%	\end{minipage}
\end{figure}
%	\hfill
%	\begin{minipage}{0.49\textwidth}
%		\centering
%		\includegraphics[width=\textwidth]{LHV_EM.pdf}
%		\caption{Syngas LHV predicted by the equilibrium model.\newline The negative impact of steam is over-estimated, and the quadratic effect of $x_{\ce{O2}}$ is milder.}
%		\label{fig:LHV_EM}
%	\end{minipage}
%\end{figure}

\renewcommand\theadalign{c}
\begin{table}[!ht]
	\centering
	\begin{tabular}{lcc@{\hskip 0.5cm} d{2.2}d{1.2}cd{2.1}}
% 	\begin{tabular}{lcc@{\hskip 0.5cm} S[table-number-alignment=center, table-figures-integer = 2, table-figures-decimal = 2]
%									  S[table-number-alignment=center, table-figures-integer = 1, table-figures-decimal = 2]
%									  c
%									  S[table-number-alignment=center, table-figures-integer = 2, table-figures-decimal = 1]}
% identical results
 \thead[l]{\;\\\nth{2} stage}  &   \thead{$x_{\ce{O2}}$\\\%}   &  \thead{SOR\\\;}   &   \multicolumn{1}{c}{\thead{LHV\\MJ/Nm$^3$}} &   \multicolumn{1}{c}{\thead{ED\\MJ/Nm$^3$}} &  \thead{AFT\\K} &  \multicolumn{1}{c}{\thead{ED/LHV\\\%}} \\
\specialrule{0pt}{2pt}{2pt}
\multirow{3}{*}{Air} & \multirow{3}{*}{21} & 0 &        4.40 &       2.31 &        1824 &           52.4 \\
                         &                     & 1 &        4.52 &       2.32 &        1829 &           51.3 \\
                         &                     & 2 &        4.47 &       2.29 &        1808 &           51.2\\
\specialrule{0pt}{2pt}{2pt}
Air-\ce{O2} & 32 & 0 &        6.14 &       2.70 &        2055 &           43.9 \\
\specialrule{0pt}{2pt}{2pt}
\multirow{3}{*}{\ce{O2}} & \multirow{3}{*}{47} & 0 &        7.26 &       2.90 &        2166 &           40.0 \\
                         &                     & 1 &        7.13 &       2.85 &        2134 &           40.0 \\
                         &                     & 2 &        6.85 &       2.78 &        2088 &           40.6 \\
\specialrule{0.5pt}{5pt}{5pt}
\ce{CO} & &  &      12.63 &        3.73 &         2663 &            29.6 \\
\ce{H2} & &  &      10.79 &        3.19 &         2524 &            29.6 \\
\ce{CH4} & &  &      35.79 &        3.40 &         2327 &             9.5 \\		
	\end{tabular}
	\captionof{table}{LHV, ED and AFT of syngas (regression model)\newline and the individual species \ce{CO}, \ce{H2} and \ce{CH4}.\label{tab:LHVEDAFT_species}}
\end{table}


%%% ER
%\begin{figure}[!ht]
%	\centering
%	\includegraphics[width=0.7\textwidth]{ER_SOR.pdf}
%	\caption{The Equivalence Ratio is lower with oxygen than with air\\
%	and slightly reduced by steam injection.\\
%	\textit{Markers: experimental points. Lines: equilibrium model (dark: constant, light: optimized)}
%\label{fig:ER}}
%\end{figure}



From the perspective of syngas applications, specific metrics can be more adequate than the LHV, such as the Energy Density (ED) and Adiabatic Flame Temperature (AFT) of the air-syngas stoichiometric mixture. The ED is proportional to the volumetric power that can be obtained when burning syngas in an internal combustion engine, for CHP applications. The AFT is relevant for processes requiring a very high flame temperature, such as glass furnaces that operate at up to 1600°C \cite{Zier2021}.

Replacing secondary air by oxygen also yields a higher ED (+26\% from 2.3 to 2.9\,MJ/Nm$^3$) and AFT (+350\,K from 1820 to 2170\,K) (Figures \ref{fig:ED} and \ref{fig:AFT} and Table \ref{tab:LHVEDAFT_species}). For a given value of LHV, both the ED and AFT are lower with steam, as steam injection slightly increases the production of \ce{CH4} (Table \ref{tab:comp}). Due to its high LHV, this small increase has a substantial effect on its share in the total syngas LHV, which rises from 6\% to 9\%. This results in a comparatively lower ED, as the ED/LHV ratio of \ce{CH4} is about 3 times lower than \ce{CO} and \ce{H2}. Due to the combined effect of a lower LHV and higher \ce{CH4} content, the ED and AFT obtained with oxygen and $\mathrm{SOR}=2$ is equivalent to the values that would be obtained using a mixture of air and oxygen such that $x_{\ce{O2}}\approx 36$\%, which corresponds to a secondary air enrichment of about 50\% \ce{O2}.


\begin{figure}[!ht]
	\begin{minipage}{0.49\textwidth}
		\centering
		\includegraphics[width=\textwidth]{ED.pdf}
		\caption{Energy Density of the stoichiometric air-syngas mixture. \textit{Regression (lines) and experimental points (markers).}}
		\label{fig:ED}
	\end{minipage}
	\hfill
	\begin{minipage}{0.49\textwidth}
		\centering
		\includegraphics[width=\textwidth]{AFT.pdf}
		\caption{Adiabatic Flame Temperature of the stoichiometric air-syngas mixture. \textit{Regression (lines) and experimental points (markers).}}
		\label{fig:AFT}
	\end{minipage}
\end{figure}

%\begin{figure}[!ht]
%	\centering
%	\includegraphics[width=0.7\textwidth]{comp_LHV.pdf}
%	\captionof{figure}{Share of \ce{CO}, \ce{H2} and \ce{CH4} in the syngas LHV.\label{fig:comp_LHV}}
%\end{figure}




\subsection{Gasifier power and efficiency}

The gasifier input power, which corresponds to the biomass flow rate, is primarily proportional to the equivalent air flow $Q_{\mathrm{air,eq}}$ (Figure \ref{fig:Pbiomass}). It also increases when replacing air by oxygen, as a result of the higher reactivity in the oxidation and reduction zones. A substantial variability of the biomass flow is observed between runs under similar conditions, and some points are even discarded as outliers. This is partly due to the substantial uncertainty on the wood chips consumption, which is not directly measured, but computed based on the duration of successive hopper loads and the average mass per hopper. The high gasifier inertia and long residence time of solid particles increase the uncertainty on the value of biomass consumption to consider, especially on relatively short runs of 1-2 hours. This uncertainty on the input power in turn affects the CGE and CCE values.

%% Power
%\begin{figure}[!ht]
%	\centering
%	\includegraphics[width=0.7\textwidth]{PbmPsg.pdf}
%	\caption{Biomass and syngas power\newline
%	increase mainly with $Q_\text{air,eq}$ and $x_{\ce{O2}}$.\newline The two semi-transparent points are outliers.\\\textit{Regression (line) and exp. points (markers)}\label{fig:PbmPsg}}
%\end{figure}
\begin{figure}[!ht]
	\centering
	\includegraphics[width=0.7\textwidth]{Pbiomass.pdf}
	\caption{Biomass consumption mainly depends on $Q_\text{air,eq}$ and $x_{\ce{O2}}$.\newline\textit{Lines: regression ($\mathrm{SOR}=0$, $\beta=2.25$). Markers: experimental points.}\newline\textit{Semi-transparent markers are outliers.}\label{fig:Pbiomass}}
\end{figure}


The actuation frequency of the ash removal rack plays an important role on the Carbon Conversion Efficiency (CCE), as shown in Figure \ref{fig:CCE}: a high frequency increases the extraction of unconverted carbon, lowering the CCE. This frequency is controlled during operation to maintain acceptable pressure losses through the reduction bed. The injection of steam appears to enhance the carbon conversion: results with $\mathrm{SOR}=2$ are associated with the lowest actuation frequency and the highest CCE. Two points have a very low CCE, probably due to too frequent ash actuations leading to a loss of carbon and an over-consumption of wood; these are considered outliers and identified by their semi-transparency in Figures \ref{fig:Pbiomass} to \ref{fig:CGE}.

A higher CCE directly translates into a higher Cold Gas Efficiency (CGE) (Figure \ref{fig:CGE}). Furthermore, the use of oxygen (full markers) yields a higher CGE than equivalent results with air (empty markers). This improvement results from the lower Equivalence Ratio with oxygen: as less biomass is burnt to sustain the process energy needs, a higher share of its LHV remains available in the syngas LHV. Combining the effects of steam on the CCE, and of steam and oxygen on the ER, the best CGE are obtained  under oxy-steam conditions: all lie between 70\% and 73\%, while most air-only, air-steam and \ce{O2}-only conditions yield a CGE between 61\% and 70\%.

%% CCE & CCE
\begin{figure}[!ht]
	\centering
	\includegraphics[width=0.6\textwidth]{CCE_ashrack.pdf}
	\caption{Carbon Conversion Efficiency is strongly influenced by the ash rack actuation frequency, and appears enhanced by steam injection.\newline\textit{Semi-transparent markers are outliers.}}
	\label{fig:CCE}
\end{figure}

\begin{figure}[!ht]
	\centering
	\includegraphics[width=0.6\textwidth]{CGE_ashrack.pdf}
	\caption{Cold Gas Efficiency is highest with a combination of oxygen and steam.\newline\textit{Semi-transparent markers are outliers.}}
	\label{fig:CGE}
\end{figure}


\subsection{Temperature in the oxidation zone}
\label{sec:res_Toxi}

The hottest zone of the gasifier is the oxidation zone, where the high temperature contributes to the efficiency of tar cracking. However, very high temperatures must be avoided as they can cause excessive mechanical wear to the concrete and stainless steel parts inside the gasifier.

Under air gasification, the average temperature in the oxidation zone $T_{\text{oxi,avg}}$ is around 980°C, with local peaks not exceeding 1020°C (Figures \ref{fig:Toxy} and \ref{fig:Toxy_ex}). The use of oxygen at the secondary stage causes a strong temperature rise, as the diluting effect of nitrogen is reduced: the average temperature increases up to 1190°C, with local peaks up to 1260°C. The fluctuations observed are due to the discontinuous transfer of char from the pyrolysis zone: when a large mass of fresh char falls to the reduction bed, the oxidation zone temperature drops as the new char suddenly cools down the second stage. The oscillations are particularly strong using oxygen, as the lower dilution of the gaseous mixture implies a lower thermal inertia.

When adding steam to oxygen, it acts as a thermal damper, reducing $T_{\text{oxi,avg}}$ to around 1040°C at $\mathrm{SOR}=2$, with local peaks not exceeding 1100°C. This temperature is similar to the one obtained with a mixture of air and oxygen at $x_{\ce{O2}}\approx 30\%$. That cooling effect of steam, mainly observed with oxygen, is due to two effects: the enthalpy required to heat it up to the local temperature, and the endothermic nature of water-gas shift and steam gasification reactions.

\begin{figure}[!ht]
	\centering
	\includegraphics[width=0.6\textwidth]{Toxy_all.pdf}
	\caption{The average temperature in the oxidation zone increases by +200°C using \ce{O2}, but this effect is compensated by the injection of steam.\newline \textit{Regression (lines) and experimental points (markers).}}
	\label{fig:Toxy}
\end{figure}

\begin{figure}[!ht]
	\centering
	\includegraphics[width=0.6\textwidth]{Toxy_ex.pdf}
	\caption{Temperature fluctuations are much stronger with oxygen, with local peaks up to 1260°C. \textit{Light dotted lines: individual thermocouples.\newline Dark solid lines: average values.}}
	\label{fig:Toxy_ex}
\end{figure}


\subsection{Light hydrocarbons}

Light hydrocarbons were measured using the on-line GC-MS-FID, during 8 of the experimental runs with air, \ce{O2} and \ce{O2}-steam. Their average concentrations in the syngas are given in Table \ref{tab:lightHCs} in the Appendix. In Figure \ref{fig:lightHCs}, these concentrations are normalized by the energy content in the syngas to avoid the bias related to the variable syngas dilution, thus expressed in mg/MJ of syngas LHV.

Amounts of ethylene and ethane are very low when no steam is injected: respectively 10\,mg/MJ (45\,mg/Nm$^3$) and 2\,mg/MJ (10\,mg/Nm$^3$). They increase when injecting steam, up to respectively 20--30\,mg/MJ and 7--13\,mg/MJ. The switch from air to oxygen does not affect their production (in mg/MJ), but their concentration increases as syngas is less diluted, respectively to 85 and 20\,mg/Nm$^3$. The concentration of benzene directly measured in the syngas is in good agreement with the values obtained by the tar protocol.

% Benzene is lowest with air only, at 30\,mg/MJ (130\,mg/Nm$^3$). Replacing secondary air by oxygen boosts it to 40--64\,mg/MJ (300-440\,mg/Nm$^3$). Adding steam further increases Benzene production in the range 50-80\,mg/MJ (380--550\,mg/Nm$^3$). 

Some methanol, propene and acetone were also found in the syngas, only under air gasification conditions (respectively 17, 2 and 6\,mg/Nm$^3$). The measurement was taken at the beginning of the stable period, and their presence might not be representative of steady-state conditions but could rather be related to the transient conditions of the gasifier.

%Il manque les résultats pour le point O260st10LTFP20; pas de point pendant la "période de stabilité" mais il en existe avant. Ils ne vont rien changer aux conclusions. A ajouter si on utilise ces résultats au final (en allant chercher des valeurs "manuellement" ?). Pour les autres, peut-être mettre aussi des valeurs qui ne sont pas pendant la "période" sélectionnée?

\begin{figure}[!ht]
	\centering
	\includegraphics[width=0.6\textwidth]{lightHCs_nobenzene.pdf}
	\caption{Ethylene and Ethane are produced in larger amounts when steam is injected.\label{fig:lightHCs}}
\end{figure}

\section{Conclusion - TO SPREAD TO DIFFERENT SECTIONS}
An experimental study has been performed on the impacts of varying amounts of air, oxygen and steam injected at the secondary stage of a pilot two-stage downdraft gasifier. The primary effect of replacing secondary air by oxygen is to increase the syngas LHV, from 4.5 up to 7.3\,MJ/Nm$^3$. This translates into a higher Energy Density (+26\%) and Adiabatic Flame Temperature (+350\,K) of the air-syngas stoichiometric mixture. This richer syngas can represent a sustainable alternative to gaseous fossil fuels for CHP applications or high-temperature industrial processes.

Using oxygen also raises the temperature in the oxidation zone of the gasifier, from 1000°C up to 1200°C, which can be harmful on the long term. Steam acts as a thermal damper, bringing the temperature back to more acceptable values with only a limited impact on the syngas LHV (-5\% to 6.9\,MJ/Nm$^3$). The steam present in the syngas condenses in the cleaning process, but it also increases the production of \ce{CO2} which results in some dilution of the dry syngas.

On the efficiency level, the injection of steam appears to enhance the conversion of carbon, which is positive on two levels: it reduces the amount of solid residue, and improves the process efficiency. The use of oxygen instead of air further increases the Cold Gas Efficiency, by enabling a lower Equivalence Ratio. The combination of oxygen and steam therefore has the potential to produce a richer syngas, with a better process efficiency.

Oxy-steam gasification has a drawback: it increases the concentration of tar in the syngas, from 60 up to 120\,mg/Nm$^3$. Class IV tars contribute to most of that increase, resulting in a higher tar dew point. However, such amounts of tars remain very low compared to most other gasifier types.

Considering their various impacts, the choice of whether to use oxygen and/or steam in two-stage downdraft gasification should be specific to each application, as trade-offs are required between their various effects and the cost of steam and oxygen production. In a future work, the energy balance of the gasifier under air, oxygen and oxy-steam conditions should be studied in the perspective of the overall process efficiency. Long-duration runs should also be performed to validate the effects observed on the conversion efficiency, based on a complete mass and energy balance. The injection of steam and oxygen at the primary stage of the gasifier should also be studied, as a next step towards an almost \ce{N2}-free syngas. Finally, the experimental results presented in this paper should be used to support the development and validation of numerical models of two-stage gasifiers.



